\section{Additional binary analysis}

\subsection{Binary constituents and their spectra}
\label{app:constituents}



The binary evaluation uses four well-known families with exact parameters and
weight spectra.  The Golay, Reed--Muller, and BCH constituents are certified in
the accompanying Lean development.  The longer quadratic-residue spectra are
standard algebraic or published enumerators and are not part of that formal
certificate.

\paragraph{Reed--Muller codes.}
The code $\operatorname{RM}(r,\ell)$ is obtained by evaluating every
multilinear polynomial in $\ell$ binary variables of total degree at most $r$
on all points of $\F_2^\ell$.  It has parameters
\[
    \left[2^\ell,\ \sum_{i=0}^r\binom{\ell}{i},\ 2^{\ell-r}\right]_2.
\]
In particular, $\operatorname{RM}(2,5)$ is a $[32,16,8]_2$ code.  Up to a
permutation of coordinates, it is also the parity extension of the primitive
triple-error-correcting BCH code $[31,16,7]_2$, and is therefore often called
the extended BCH code $[32,16,8]_2$~\cite{macwilliams}.  Here ``equivalent''
means that the two descriptions differ only by a relabeling of coordinates;
they have the same weight spectrum and give exactly the same BAA transition.
Thus the $\operatorname{RM}(2,5)$ calculation below is already a rate-$1/2$
BCH result.

\paragraph{BCH codes.}
A primitive binary BCH code is a cyclic code of length $2^\ell-1$ specified by
requiring its code polynomials to vanish at a consecutive set of powers of a
primitive element of $\F_{2^\ell}$.  The chosen roots give a designed-distance
lower bound; extending the code by an overall parity coordinate can increase
its length and distance.  A subcode is simply a linear subspace of the
parent code.  The rate-$1/4$ block used here is an explicitly generated
$[32,8,12]_2$ subcode of an extended primitive $[32,11,12]_2$ BCH code.
Appendix~\ref{app:binary-generator} gives its generator and its exact exhaustive
weight enumeration, so the analysis does not rely only on the BCH designed
distance.

\paragraph{The extended Golay code.}
The extended binary Golay code is the exceptional self-dual $[24,12,8]_2$
code.  Along with the extended BCH code (equivalently
$\operatorname{RM}(2,5)$), it supplies a compact rate-$1/2$ block with a known
complete weight enumerator.

\paragraph{Extended quadratic-residue codes.}
Extending the binary quadratic-residue code of prime length $p\equiv-1\pmod 8$
gives a doubly even self-dual code of length $p+1$ and rate $1/2$.  We evaluate
the concrete constituents $[48,24,12]$, $[80,40,16]$, $[104,52,20]$, and
$[168,84,24]$.  For the first three, Type-II extremality and Gleason's theorem
fix the complete weight enumerator from the displayed parameters.  For length
$168$, we use the exact weight distribution determined by Su et
al.~\cite{su08}; Appendix~\ref{app:qr-spectrum} records every nonzero
coefficient used in our calculation.




\subsection{Accumulator transition proof}
\label{app:binary-analysis}
The run-counting argument for Lemma~\ref{lem:binary-transition} is recalled here for completeness.
\begin{proof}
A nonzero input coordinate toggles the accumulator state.  Decomposing the
output into alternating zero-runs and one-runs gives a bijection with two run
compositions.  The positive lengths of the $r_w$ one-runs contribute the first
binomial coefficient.  The $s_w$ required zero gaps contribute the second.
The inverse difference map proves both injectivity and surjectivity of this
correspondence.  Invertibility excludes a zero output for nonzero input.
\end{proof}



\subsection{Binary asymptotic calculations}
\label{app:binary-asymptotics}

Fix $w$ and a target $0<\delta<1/2$.  The exact transition~\eqref{eq:B3} gives, for
fixed $h$,
\[
    P_n^{\rm bin}(w,h)
      =
    \Theta_{w,h}\!\left(n^{-\lceil w/2\rceil}\right).
\]
The second accumulator tail from fixed $h$ converges to a nonzero beta-tail
constant.  Splitting the intermediate-weight sum into bounded and growing
ranges therefore gives
\[
    \boxed{
    \sum_h P_n^{\rm bin}(w,h)Q_n^{\rm bin}(h,\lfloor\delta n\rfloor)
      =
    \Theta_{w,\delta}\!\left(n^{-\lceil w/2\rceil}\right).
    }
    \eqnumlabel{eq:B6}
\]
Hence the one-active-block contribution is
\[
    \frac{k}{\sigma}
    \sum_{w=1}^{m}A_w^{(0)}
    \Theta_{w,\delta}\!\left(n^{-\lceil w/2\rceil}\right).
    \eqnumlabel{eq:B7}
\]
If the local distance is $d$, this is at most
\[
    O\!\left(n^{1-\lceil d/2\rceil}\right),
\]
with the exact leading term determined by all local weights attaining the
smallest value of $\lceil w/2\rceil$.

Using only local distance, a fixed family touching $t$ blocks has the coarse
scale
\[
    O\!\left(n^{t-\lceil td/2\rceil}\right).
    \eqnumlabel{eq:B8}
\]
Thus a constant local distance of at least three is consistent with a vanishing
fixed-$t$ sparse tail.  The earlier binary BA proof assumes $d>4$ because it
uses coarser inequalities.  For a fixed local block, its exact spectrum and the
two-accumulator kernel~\eqref{eq:B4} give the sharper bound.

\paragraph{Linear weights.}

We use the entropy analysis of accumulator transitions developed for RAA/RMA
codes~\cite[Sections~III--IV]{kliewer08}.  For the outer stage, the elementary
coefficient bound $[z^w]W_0(z)^{n/m}\le W_0(z)^{n/m}z^{-w}$ for $z>0$
gives the direct-sum spectrum the upper exponent
\[
    \boxed{
    \mathcal A_{C_0}(\alpha)
      :=
    \inf_{z>0}
    \left\{
      \frac1m\log W_0(z)-\alpha\log z
    \right\}.
    }
    \eqnumlabel{eq:B9}
\]
Indeed,
\[
    A_{\lfloor\alpha n\rfloor}^{\rm out}
      \le
    \exp\!\left(n\mathcal A_{C_0}(\alpha)+O(\log n)\right).
\]
For linear input and output weights $u n$ and $v n$, the binary accumulator
transition exponent is
\[
    \boxed{
    \Gamma_{\rm bin}(u,v)
      :=
    vH\!\left(\frac{u}{2v}\right)
    +(1-v)H\!\left(\frac{u}{2(1-v)}\right)
    -H(u),
    }
    \eqnumlabel{eq:B10}
\]
with feasibility $u/2\le v\le1-u/2$.  For $\delta<1/2$, this exponent is
increasing in the final output fraction up to $1/2$, so the low-weight tail is
maximized at the boundary $v=\delta$.

The binary BAA first-moment exponent is therefore
\[
    \boxed{
    \Phi^{\rm bin}_{C_0,\delta}(\alpha,\eta)
      =
    \mathcal A_{C_0}(\alpha)
    +\Gamma_{\rm bin}(\alpha,\eta)
    +\Gamma_{\rm bin}(\eta,\delta).
    }
    \eqnumlabel{eq:B11}
\]
For a fixed $\varepsilon>0$, a sufficient condition for the contribution
from outer weights $w\ge\varepsilon n$ to vanish is
\[
    \sup_{\substack{\alpha\ge\varepsilon,\,\eta\\\text{feasible}}}
       \Phi^{\rm bin}_{C_0,\delta}(\alpha,\eta)<0
    \eqnumlabel{eq:B12}
\]
on this restricted domain.  The unrestricted supremum is at least zero
because the exponent tends to zero at the sparse boundary.
For a fixed block $C_0$, \eqref{eq:B4} gives the exact
finite-length first moment, \eqref{eq:B7} controls fixed-block supports, and \eqref{eq:B12}
controls linear weights.




\subsection{Binary numerical certificates}
\label{app:binary-certificate}

\paragraph{Rate $1/2$: extended Golay.}
The extended binary Golay code has parameters $[24,12,8]$ and spectrum
\[
    W_{\rm Golay}(z)
      =
    1+759z^8+2576z^{12}+759z^{16}+z^{24}.
\]
The numerical BAA transition is
\[
    \boxed{\delta^\star_{\rm Golay}\approx0.105.}
\]
The contribution from one active block has order $n^{-3}$.

\paragraph{Rate $1/2$: extended BCH / $\operatorname{RM}(2,5)$.}
The extended BCH code $[32,16,8]$, equivalently the second-order Reed--Muller
code $\operatorname{RM}(2,5)$, has spectrum
\[
\begin{aligned}
    W_{\operatorname{RM}(2,5)}(z)
      ={}&
    1+620z^8+13888z^{12}+36518z^{16}
    \\
    &+13888z^{20}+620z^{24}+z^{32}.
\end{aligned}
\]
The numerical BAA transition is
\[
    \boxed{\delta^\star_{\mathrm{eBCH}/\operatorname{RM}(2,5)}
      \approx0.106.}
\]
This is our compact rate-$1/2$ BCH baseline.  Among the two compact implemented
blocks, it gives the larger transition: about $96.5\%$ of the binary GV
 relative distance.  The longer rate-$1/2$ quadratic-residue spectra below
 improve on this value.  The one-block contribution again has order $n^{-3}$.

\paragraph{Finite-length $40$-bit certificate for the implemented blocks.}
We evaluate the complete sum~\eqref{eq:B4} at the benchmark length
$n=2{,}097{,}024$ and $L=\lfloor0.10n\rfloor=209{,}702$.  To make the
calculation both exhaustive and reproducible, first observe that
$Q_n^{\rm bin}(h,L)$ is the upper tail
$\Prb[X\ge x_h]$ for
$X\sim\operatorname{Hypergeom}(n,L,h)$ and $x_h=\lceil h/2\rceil$.  If
\[
 \rho_h=
 \frac{L-x_h}{x_h+1}
 \frac{h-x_h}{n-L-h+x_h+1},
\]
then successive terms of this tail have decreasing ratio, and hence
\[
 Q_n^{\rm bin}(h,L)
 \le
 \frac{\binom{L}{x_h}\binom{n-L}{h-x_h}}
      {\binom{n}{h}(1-\rho_h)}.
 \eqnumlabel{eq:binary-hypergeom-tail}
\]
We compute the coefficients $A_w^{\rm out}$ exactly through $w=32$ and sum
the resulting bound over every feasible intermediate weight $h$.  For the
remaining weights, the coefficient inequality
\[
 A_w^{\rm out}
 \le
 \exp\!\left(B\log W_0(e^\theta)-w\theta\right),
 \qquad B=n/m,
 \eqnumlabel{eq:binary-coefficient-bound}
\]
 holds for every real $\theta$.  We minimize this bound numerically.  An
 explicit consecutive-ratio check shows that the bounded transition summand is
 unimodal on each of the even and odd $h$ subsequences; we maximize the two
 subsequences separately and multiply their maximum by the number of feasible
 $h$.  No $w>4L$ term can contribute, because each binary
accumulator has output weight at least half its nonzero input weight.  A full
one-bit outward-rounding allowance is applied independently to the $w\le32$
and $w>32$ subtotals.

The resulting complete first-moment certificates are
\[
 \boxed{
 \begin{aligned}
  \mu_{\rm bin}(n,L;C_{\rm Golay})
    &\le 2^{-46.386},\\
  \mu_{\rm bin}(n,L;\operatorname{RM}(2,5))
    &\le 2^{-47.093}.
 \end{aligned}}
 \eqnumlabel{eq:binary-40-certificate}
\]
For comparison, the $w>32$ subtotal is at most $2^{-155.759}$ for Golay and
$2^{-158.897}$ for Reed--Muller; the complete bound is controlled by $w=8$.
 Repeating the calculation with cutoffs $24$ and $32$ gives the same displayed
 complete bounds.  The binary first-moment script in the accompanying
 repository~\cite{chosenbaaartifact} reproduces the calculation.
 By~\eqref{eq:B5}, both implemented chosen-block ensembles
therefore exceed a $40$-bit one-time sampling-failure target at relative
distance $0.10$.  This certifies the chosen rows themselves; it does not turn
the original $\sigma=32$ row, documented at $20$ bits, into a matched $40$-bit
baseline.




\section{Additional large-field analysis}

\subsection{Large-field asymptotic calculations}
\label{app:field-asymptotics}

\paragraph{Sparse supports.}

Suppose an admissible global support touches $t$ local blocks.
In the no-surplus regime,
\[
    a+b\le r-1.
\]
The first accumulator output has at least $s-a$ nonzero-state gaps, and after
the second stage any final output of weight $j$ satisfies
\[
    j\ge s-a-b.
\]
Using the full projective cancellation budget gives
\[
    s-a-b\ge s-r+1\ge t(d-1)+1.
\]

A direct fixed-$t$ asymptotic analysis of the exact kernel gives the scale
\[
    \boxed{
    P_t
      =
    O\!\left(
      n^{-t(d-2)-1}
    \right)
    }
    \eqnumlabel{eq:6}
\]
when $q$ is sufficiently large that surplus cancellations are negligible.

In particular,
\[
    P_1=O(n^{1-d}).
\]
The canonical one-block event gives a matching lower bound of the same
polynomial order, so the natural sparse tail is
\[
    \boxed{
    \Prb[\text{sparse bad event}]
      =
    \Theta(n^{1-d}).
    }
    \eqnumlabel{eq:7}
\]

For example:
\[
\begin{array}{c|c}
d & \text{sparse tail scale}\\
\hline
2 & n^{-1}\\
3 & n^{-2}\\
5 & n^{-4}\\
7 & n^{-6}\\
9 & n^{-8}
\end{array}
\]

For $q-1\gtrsim n$, a surplus cancellation can recover at most one geometric
power of $n$ while costing $1/(q-1)$, so surplus layers remain controlled.
For $q\approx2^{128}$ and the lengths considered here, the resulting surplus
terms are negligible on the reported scale.

\paragraph{Linear weights.}

Let $C_0$ be an $[m,\sigma,d]$ MDS code, with
\[
    d=m-\sigma+1,
    \qquad R=\sigma/m=k/n.
\]
Its local support--dimension polynomial is
\[
    \boxed{
    F_{m,\sigma}(z,u)
      =
    1+
    \sum_{\ell=d}^{m}
    \binom{m}{\ell}z^\ell u^{\ell-d+1}.
    }
    \eqnumlabel{eq:L1}
\]
Let $p_0,p_d,\ldots,p_m$ be the fractions of local blocks in each exact
support class, and put
\[
    \lambda:=\sum_{\ell=d}^{m}p_\ell.
\]
Define the normalized global support and shortening dimension by
\[
    \alpha
      :=
    \frac1m\sum_{\ell=d}^{m}\ell p_\ell,
    \qquad
    \beta
      :=
    \frac1m\sum_{\ell=d}^{m}(\ell-d+1)p_\ell.
\]
Since $d-1=m-\sigma=m(1-R)$,
\[
    \boxed{
    \beta=\alpha-(1-R)\lambda.
    }
    \eqnumlabel{eq:L2}
\]
The support-pattern entropy per output coordinate is
\[
    \boxed{
    \cS_{m,\sigma}(p)
      :=
    \frac1m
    \left[
      H(p_0,p_d,\ldots,p_m)
      +
      \sum_{\ell=d}^{m}p_\ell\log\binom{m}{\ell}
    \right].
    }
    \eqnumlabel{eq:L3}
\]
Here $H$ uses natural logarithms unless a quantity is explicitly divided by
$\log 2$.

Normalize the first- and second-stage zero-return counts and intermediate
support by
\[
    x:=a/n,
    \qquad
    y:=b/n,
    \qquad
    \eta:=h/n.
\]
The first accumulator contributes
\[
\boxed{
\begin{aligned}
\Gamma_1(\alpha,\eta,x)
={}&
(\alpha-x)
H\!\left(\frac{x}{\alpha-x}\right)
+
(1-\eta)
H\!\left(\frac{x}{1-\eta}\right)
\\
&+
\eta
H\!\left(\frac{\alpha-x}{\eta}\right)
-H(\alpha),
\end{aligned}
}
\eqnumlabel{eq:L4}
\]
and, at final relative weight $\zeta$, the second contributes
\[
\boxed{
\begin{aligned}
\Gamma_2(\eta,\zeta,y)
={}&
(\eta-y)
H\!\left(\frac{y}{\eta-y}\right)
+
(1-\zeta)
H\!\left(\frac{y}{1-\zeta}\right)
\\
&+
\zeta
H\!\left(\frac{\eta-y}{\zeta}\right)
-H(\eta).
\end{aligned}
}
\eqnumlabel{eq:L5}
\]
A direct derivative calculation gives the unconstrained maximizing final weight
\[
    \boxed{
    \zeta_\star=1-\frac{y}{\eta}.
    }
    \eqnumlabel{eq:L6}
\]
Therefore define
\[
    \widehat\Gamma_2(\eta,\delta,y)
      :=
    \Gamma_2\!\left(
      \eta,
      \min\left\{\delta,1-\frac{y}{\eta}\right\},
      y
    \right).
\]
This distinction matters for target distances exceeding $1/2$.

In the large-field no-surplus region,
\[
    x+y\le\beta,
\]
and the complete support/trace first-moment exponent is
\[
    \boxed{
    \Phi_{m,\sigma,\delta}
      =
    \cS_{m,\sigma}(p)
      +
    \Gamma_1(\alpha,\eta,x)
      +
    \widehat\Gamma_2(\eta,\delta,y).
    }
    \eqnumlabel{eq:L7}
\]
For finite $q$, the corresponding exponent includes
\[
    \min\left\{
      0,
      \beta\log q-(x+y)\log(q-1)
    \right\}.
\]
The exponent tends to zero as the active-block fraction tends to zero, so
its unrestricted supremum cannot be negative.  The numerical transition
estimates below track the zero crossing of a nonzero stationary branch of
\eqref{eq:L7}.  We write $\Phi^{\mathrm{int}}_{m,\sigma,\delta}$ for the
reported value on that branch.  A negative branch value alone does not
certify a uniform bound.  The finite calculation in
Appendix~\ref{app:rs-certificate} instead bounds an explicitly specified
domain away from the sparse boundary and covers the complementary supports
separately.




\subsection{Worked Reed--Solomon evaluations}
\label{app:rs-evaluation}

\paragraph{Rate $1/2$.}

Take
\[
    C_0=\operatorname{RS}[16,8,9].
\]
Then
\[
    F_{16,8}(z,u)
      =
    1+
    \sum_{\ell=9}^{16}
    \binom{16}{\ell}z^\ell u^{\ell-8},
\]
and
\[
    \beta=\alpha-\frac12\lambda.
\]
Numerical optimization of the support/trace exponent gives the transition
\[
    \boxed{
    \delta^\star_{16,8}\approx0.305.
    }
    \eqnumlabel{eq:R1}
\]
At the transition, a representative saddle is approximately
\[
\begin{aligned}
    \lambda &\approx0.500, &
    \alpha &\approx0.448, &
    \beta &\approx0.198,\\
    x &\approx0.108, &
    y &\approx0.090, &
    \eta &\approx0.367.
\end{aligned}
\]
At this saddle, the available projective dimension is almost exactly exhausted
by the zero-return constraints:
\[
    \boxed{x+y\approx\beta.}
\]
At $\delta=0.30$, the reported interior-branch exponent is approximately
\[
    \Phi_{16,8,0.30}^{\mathrm{int}}/\log 2\approx-0.017
\]
bits per output coordinate.

The local distance is $d=9$, so the sparse one-block scale is $n^{-8}$.  At the
primary dimension $k=2^{20}$ and output length $n=2^{21}$, with $\delta=0.30$,
the marked one-block term is approximately $2^{-137.57}$.  That term and
the negative linear-weight exponent alone do not control every support
size.  Theorem~\ref{thm:rs-finite} supplies the complete finite-length bound
at this operating point.

\paragraph{Rate $1/4$.}

Take
\[
    C_0=\operatorname{RS}[16,4,13].
\]
The local support--dimension polynomial is
\[
    \boxed{
    F_{16,4}(z,u)
      =
    1+560z^{13}u+120z^{14}u^2+16z^{15}u^3+z^{16}u^4.
    }
    \eqnumlabel{eq:R3}
\]
With
\[
    \lambda=p_{13}+p_{14}+p_{15}+p_{16},
\]
equation~\eqref{eq:L2} becomes
\[
    \boxed{
    \beta=\alpha-\frac34\lambda.
    }
    \eqnumlabel{eq:R4}
\]
The larger redundancy sharply reduces the projective cancellation budget.

Numerical optimization of~\eqref{eq:L7} gives the estimated transition
\[
    \boxed{
    \delta^\star_{16,4}\approx0.604.
    }
    \eqnumlabel{eq:R5}
\]
At the transition, a representative saddle is approximately
\[
\begin{aligned}
    \lambda &\approx0.670, &
    \alpha &\approx0.660, &
    \beta &\approx0.158,\\
    x &\approx0.065, &
    y &\approx0.093, &
    \eta &\approx0.666.
\end{aligned}
\]
Here too, the zero-return constraints nearly exhaust the projective dimension:
\[
    \boxed{x+y\approx\beta.}
\]
At this saddle,
\[
    1-\frac{y}{\eta}\approx0.861>\delta^\star_{16,4},
\]
so the final-weight constraint is active and the maximizing value is
$\zeta=\delta$.

Conditional on a block being active, the optimizing local-support distribution
is approximately
\[
\begin{array}{c|c}
\text{local support} & \text{conditional fraction}\\
\hline
13 & 0.002\\
14 & 0.022\\
15 & 0.190\\
16 & 0.787
\end{array}
\]
Thus the transition is dominated by dense local shortenings rather than
minimum-weight words.

In bits per output coordinate, the exponent decomposes approximately as
\[
    \cS_{16,4}\approx+0.132,
    \qquad
    \Gamma_1\approx-0.068,
    \qquad
    \Gamma_2\approx-0.064,
\]
which sum to zero at the transition.  This is about $81.3\%$ of the
large-field rate-$1/4$ GV relative distance.

\paragraph{Sparse tail and finite-length margins.}

Here $d=13$, so the one-active-block structural tail has scale
\[
    \boxed{n^{-12}.}
\]
The leading marked term is
\[
    \boxed{
    \mu_1^{\rm mark}
      \sim
    \frac{13!}{16}
    \binom{16}{13}^{\!2}
    \left(\frac{\delta}{1-\delta}\right)^{13}
    n^{-12}.
    }
    \eqnumlabel{eq:R6}
\]
At dimension $k=2^{20}$ and output length $n=2^{22}$, the numerical estimates
are
\[
\begin{array}{c|c}
\delta & -\log_2\mu_1^{\rm mark}\\
\hline
0.550 & 213.44\\
0.580 & 211.15\\
0.600 & 209.60\\
0.604 & 209.26
\end{array}
\]
The estimated interior-branch margins at this length are
\[
\begin{array}{c|c|c}
\delta
& \Phi^{\mathrm{int}}_{16,4,\delta}/\log 2
& -n\Phi^{\mathrm{int}}_{16,4,\delta}/\log 2\text{ (bits)}\\
\hline
0.600 & -0.017 & 70828\\
\text{near }0.604 &
[-0.002,-5.15\mathbin{\times}10^{-5}] & 216\text{--}6708
\end{array}
\]
For $k=2^{20}$ and $n=2^{22}$, the high-precision point used in the
first-moment calculation rounds to
\[
    \boxed{
    C_0=\operatorname{RS}[16,4,13],
    \qquad
    \delta=0.604
    }
    \eqnumlabel{eq:R7}
\]
gives a numerical margin of roughly $160$ bits.  This does not certify that
$d_{\min}>0.604n$.  At $\delta=0.60$, the estimated margin is much larger.




\subsection{Repeat--accumulate comparison bounds}
\label{app:ra-bound}

For a full finite-length comparison, put $h_0=rw$.  Applying
Lemma~\ref{lem:nonbinary-placement} stage by stage and
Lemma~\ref{lem:projective-cancellation} to the $w$-dimensional shortening of
$w$ active repetition blocks gives
\[
\boxed{
\begin{aligned}
 \Prb[d_{\min}(\cC_{\rm RA^t})\le L]
 &\le \mu_{\rm RA^t}(N,L;r)\\
 &:={}
 \sum_{w=1}^{k_0}\binom{k_0}{w}
 \sum_{\substack{a_1,\ldots,a_t\ge0\\
                   h_1,\ldots,h_t\ge1\\ h_t\le L}}
 \chi_q\!\left(w,\sum_{i=1}^t a_i\right)
 \prod_{i=1}^t \cG_N(h_{i-1},h_i,a_i).
\end{aligned}}
\eqnumlabel{eq:rat-complete-bound}
\]
Infeasible trace tuples contribute zero through $\cG_N$.  The outer binomial
coefficient counts the possible active message supports, and their shortening
dimension is $w$.  Thus this is the complete support/trace union bound, not only
the weight-one no-return contribution.  That leading contribution still
explains the sparse orders $N^{1-r}$ for RA$^2$ and
$N^{1-r(t-1)}$ for fixed $r,t$.




\section{Structural-container calculations}

\subsection{Counting structural containers}
\label{app:container-count}



For fixed permutations, define
\[
    Z_t
      :=
    \bigl|\{T\subseteq[n]:|T|=t,\ \kappa(U(T))<t\}\bigr|.
\]
By \eqref{eq:G5}, every set counted by $Z_t$ supports a nonzero BAA codeword for every
choice of diagonal coefficients.

Endpoint membership changes the size of a difference neighborhood.  For
$e\in\{0,1\}$, let $R_n^e(t,r)$ count the $t$-subsets with $r$ runs and with
$e=1$ exactly when $n$ belongs to the set.  Then
\[
\begin{aligned}
    R_n^0(t,r)
      &:=\binom{t-1}{r-1}\binom{n-t}{r},\\
    R_n^1(t,r)
      &:=\binom{t-1}{r-1}\binom{n-t}{r-1},\\
    R_n(t,r)
      &:=R_n^0(t,r)+R_n^1(t,r)
       =\binom{t-1}{r-1}\binom{n-t+1}{r}.
\end{aligned}
\eqnumlabel{eq:G6}
\]
For a set of type $(t,r,e)$, the neighborhood has exact size
\[
    |\partial T|=t+r-e.
\]

For the local block-rank count define
\[
    B_{m,h}(x,y)
      :=
    \sum_{\ell=0}^{m}\binom m\ell x^\ell y^{\min(\ell,h)},
    \eqnumlabel{eq:G7}
\]
and
\[
    A_{n,m,h}(u,c)
      :=
    \sum_{z\le c}[x^uy^z]B_{m,h}(x,y)^{n/m}.
    \eqnumlabel{eq:G8}
\]

\begin{theorem}[Exact structural-container expectation]
Assume $m\mid n$, fix $1\le t\le n$, and sample independent
$P_1,P_2\gets\operatorname{Sym}([n])$.  For
$r_1,r_2\ge1$ and $e_1,e_2\in\{0,1\}$, set
\[
    u_1:=t+r_1-e_1,
    \qquad
    u_2:=u_1+r_2-e_2.
\]
Let $\mathcal I_t$ contain the tuples $(r_1,r_2,e_1,e_2)$ that satisfy
$r_1\le t$, $r_2\le u_1$, and $u_1,u_2\le n$.
Then the expected number of direct structural containers $T$ of size $t$ is
\[
\boxed{
\begin{aligned}
\E Z_t
=
\sum_{(r_1,r_2,e_1,e_2)\in\mathcal I_t}
&R_n^{e_1}(t,r_1)
\frac{R_n^{e_2}(u_1,r_2)}{\binom n{u_1}}
\\[-1mm]
&\cdot
\frac{A_{n,m,h}(u_2,t-1)}
     {\binom n{u_2}}.
\end{aligned}
}
\eqnumlabel{eq:G9}
\]
\end{theorem}

\begin{proof}
Partition the $t$-subsets by their first run count and endpoint flag.  Equation
\eqref{eq:G6} gives $R_n^{e_1}(t,r_1)$ choices, and the first neighborhood has size
$u_1$.  The permutation $P_2$ maps this neighborhood to a uniform
$u_1$-subset.  A second application of \eqref{eq:G6} gives the ratio
$R_n^{e_2}(u_1,r_2)/\binom n{u_1}$ and a neighborhood of size $u_2$.  The
permutation $P_1$ maps that neighborhood to a uniform $u_2$-subset.  Equation
\eqref{eq:G8}, divided by $\binom n{u_2}$, is exactly the probability that its local
block rank is at most $t-1$.  Multiplying the factors and summing the disjoint
cases proves \eqref{eq:G9}.
\end{proof}

Impossible parameter choices contribute zero under the usual binomial
convention.  Equation~\eqref{eq:G9} counts each final support once and contains no field
values or zero-return choices.  Omitting $e_1,e_2$ changes neighborhood sizes
by at most two and gives the same entropy exponent, but not the same finite
expectation.




\subsection{Structural optimization details}
\label{app:structural-optimization}

The rate-$1/2$ structural saddle has normalized statistics
\[
\begin{aligned}
|T|/n&\approx0.308,& r_1/n&\approx0.101,\\
|\partial T|/n&\approx0.409,& r_2/n&\approx0.120,\\
|U(T)|/n&\approx0.530,& \kappa(U(T))/n&\approx0.308.
\end{aligned}
\eqnumlabel{eq:G11}
\]
At rate $1/4$ the corresponding values are
\[
\begin{aligned}
|T|/n&\approx0.609,& r_1/n&\approx0.103,\\
|\partial T|/n&\approx0.712,& r_2/n&\approx0.070,\\
|U(T)|/n&\approx0.782,& \kappa(U(T))/n&\approx0.609.
\end{aligned}
\eqnumlabel{eq:G12}
\]
Moderate-length integer-program experiments on the exact min-cut formulation
produce block-saturated witnesses consistent with these profiles.  These
experiments and all numbers in \eqref{eq:G10}--\eqref{eq:G12} are diagnostic rather than theorem
statements.




\section{Implementation and benchmark details}

\subsection{Benchmark reproduction details}
\label{app:benchmarks}

\paragraph{Controlled randomness and timed region.}
All three implementations use $n=2{,}097{,}024$, the largest multiple of
$\operatorname{lcm}(24,32,64)=192$ below $2^{21}$, and
$k=n/2=1{,}048{,}512$.  Divisibility by $24$, $32$, and $64$ lets the Golay,
Reed--Muller, and original random-block constituents tile one common input
without truncation.  All rows therefore use the same input, two accumulator
layers, two independently seeded pseudorandom permutations, one CPU thread, and
timed $n$-to-$k$ operation.  The input and permutation seeds are generated
before timing from the same deterministic seeds and in the same order.

The original-BAA-style control uses two prefix-sum/permutation layers followed by a
dense random systematic local map.  The two permutations are sampled
independently, and the timed computation evaluates the $n$-to-$k$
multiplication $x\leftarrow x+Ry$.  Chosen-block BAA replaces the random local
matrices with the chosen constituent and hence needs no local-matrix stream.

Input generation, permutation-seed derivation, and workspace allocation occur before
timing.  Both accumulator/permutation layers and the local outer multiplication
remain inside it.  The original baseline generates its random local matrices
during that operation, whereas the chosen constructions evaluate their
constituent-code circuits.  The benchmark supports the common non-power-of-two
length without a stored $n$-entry permutation table.  The mathematical ensemble
analysis continues to sample uniform permutations.

For the large-field configuration menu, every row fixes $k=2^{20}$,
$\F_{2^{128}}$, one CPU thread, and the same encode-only boundary, while rate,
length, local map or repetition, and accumulator count vary intentionally.  The
rate-$1/2$ rows have $n=2^{21}$; repetition-$3$ RA$^2$ has
$n=3\cdot2^{20}$; and RS $[16,4,13]$ BAA has rate $1/4$ and $n=2^{22}$.
The reported target distances are $0.10$ for repetition-$2$ RA$^2$, $0.30$
for rate-$1/2$ BAA, $0.40$ for repetition-$2$ RA$^3$ and repetition-$3$
RA$^2$, and $0.60$ for rate-$1/4$ BAA.  Each lies below that configuration's
estimated trace transition; lower-rate configurations are evaluated at larger
distances rather than at the rate-$1/2$ reference $0.10$.
Each stage has independent diagonal weighters and interleavers.  In every case,
permutation and diagonal generation, input generation, and workspace allocation
are outside the timed region; every full scan--multiply--scatter layer and the
outer transpose multiplication or repetition reduction are inside it.

\paragraph{Platform and procedure.}
On an Intel Core i9-14900HX with $32$ logical processors, we compiled libOTe in
Release mode with MSVC~14.44 and C++20 on Windows NT~10.0.26200, with SSE and
 PCLMUL enabled.  We exclude setup, input generation, permutation construction,
 and workspace allocation; all scan/permutation layers and the outer
 multiplication are timed.  The headline binary construction campaign uses
exactly one CPU thread and binds each process to logical processor~2 (Windows
affinity mask \texttt{0x4}).  We report the median of six process medians, each
over $31$ calls after warm-up.  Across the six processes, each of the three
implementations occupies each benchmark-order position twice.  The separate,
unpinned two-mode scaling campaigns reverse order halfway.  Turbo and thermal
state are not controlled; results are host-specific.
Each large-field row likewise uses six $31$-call process medians after warm-up;
paired BAA and RA$^t$ campaigns alternate which implementation runs first.




\subsection{CPU control and permutation ablations}
\label{app:ablations}

For context, the original BAA paper reports $21.4$~ms for its aggressive
$\sigma=8$ row and $30.2$~ms for its conservative $\sigma=32$ row at
$n=2^{21}$ on an Intel Core i7-13700H~\cite{baa26}.  Table~\ref{tab:cpu-chosen}
instead uses $n=2{,}097{,}024$, the largest multiple of
$\operatorname{lcm}(24,32,64)$ below $2^{21}$, so that the Golay,
Reed--Muller, and random-block constituents all tile the same input without a
shortened final block.  This common-workload choice is $128$ symbols below the
original length.  In a separate affinity-controlled path ablation, the
corrected $\sigma=32$ control took $31.895$~ms at the common length and
$31.992$~ms at $n=2^{21}$ under the same on-the-fly permutation implementation;
their six-process ranges, $29.147$--$32.997$ and $28.999$--$33.039$~ms, overlap.
A precomputed-permutation variant at the common length took $33.688$~ms in the
same ablation.  Thus the $128$-symbol length reduction has no measurable
penalty once the permutation implementation is held fixed.
Table~\ref{tab:cpu-chosen}'s $29.596$~ms
control comes from the separate balanced three-way campaign and is used only
for its within-campaign ratios.  The original paper's executable, compiler
environment, and processor (i7-13700H rather than i9-14900HX) also differ from
ours, so hardware labels alone do not make the absolute timings directly
comparable.




\subsection{Projected PCG throughput}
\label{app:pcg-throughput}

\paragraph{Estimated single-thread PCG throughput.}
The original BAA evaluation reports about $8$~ms of CPU expansion work per
batch, in addition to the linear-code encoding, and excludes an approximately
$10$~ms one-time seed-tree cost from its throughput figure because that cost is
amortized across later invocations~\cite{baa26}.  Following the same convention,
we combine its recurring $8$~ms estimate with our one-thread Golay encoding time:
\[
  \frac{1{,}048{,}512}{(22.810+8)\ \mathrm{ms}}
  \approx 34.0\ \text{million OTs/s}.
\]
This projected rate excludes the one-time $10$~ms cost, as does the original
paper; it is not an integrated measurement.  We do not extrapolate an
end-to-end multithreaded rate because the corresponding parallel expansion cost
was not measured; the encoding-only scaling results remain below.




\section{Formalization and proof-audit details}

\subsection{Formalization details}
\label{app:formalization}

\subsubsection{Representation and proof architecture}

A length-$n$ word over $F$ is represented as a function
$\operatorname{Fin}(n)\to F$, and a linear code is a submodule of that vector
space.  Hamming weight uses mathlib's \texttt{hammingNorm}.  The accumulator,
difference, permutation, diagonal, and monomial transformations are linear
equivalences.  The repeated outer code uses an explicit equivalence between
$\operatorname{Fin}(n)$ and
$\operatorname{Fin}(n/m)\times\operatorname{Fin}(m)$, which removes hidden
divisibility and coordinate-order assumptions.

Ensemble probabilities are represented as ratios of finite cardinalities in
$\mathbb Q$.  Both first-moment arguments are therefore exact finite counts;
no measure-theoretic layer is needed.  Generic rank is defined over the fraction
field of a multivariate polynomial ring.  The two diagonal layers use disjoint
variable families, so different Cauchy--Binet selections yield different
monomials.

The concrete binary codes also have a computation-oriented representation based
on fixed-width bit vectors.  Conversion lemmas connect these evaluators to the
abstract submodules.  Consequently, a computed spectrum certifies the code in
the corresponding mathematical theorem.

\subsubsection{Verified statements}

The development contains complete proofs of the following finite claims.

\begin{itemize}
    \item The BAA transformations are invertible and preserve dimension.  The
          block-direct-sum enumerator and the binary identities
          \eqref{eq:B1}--\eqref{eq:B5} are
          exact finite theorems.
    \item The extended Golay and $\operatorname{RM}(2,5)$ constructions have
          the dimensions, minimum distances, and spectra stated above.  The
          eight-row $[32,8,12]$ code has the stated spectrum and lies in a
          concretely defined extended primitive BCH code.
    \item Reed--Solomon evaluation codes have the claimed dimension and
          minimum distance.  Their MDS shortening dimensions give
          \eqref{eq:5}, \eqref{eq:L1}, \eqref{eq:L2}, \eqref{eq:R3}, and
          \eqref{eq:R4}.
    \item The zero-return placement count \eqref{eq:1}, projective bound
          \eqref{eq:2}, master bound \eqref{eq:3}, and corrected support-excess
          bound \eqref{eq:4} are formal
          consequences of finite counting and linear algebra.
    \item The backward parity-check identities \eqref{eq:G1}--\eqref{eq:G5}, run and block-rank
          counts \eqref{eq:G6}--\eqref{eq:G8}, and exact structural-container
          expectation \eqref{eq:G9},
          including the endpoint cases, are proved.  The generic rank equals
          the maximum size of a two-layer
          vertex-disjoint path packing with the local MDS capacities.
\end{itemize}

The quadratic-residue rows in Table~\ref{tab:binary-qr} use exact external
spectra and are deliberately outside this list of Lean-verified statements.

The path-packing result is stronger than the neighborhood upper bound used in
\eqref{eq:G5}.  Its proof first establishes a rectangular Cauchy--Binet identity.  It
then identifies nonzero minors of the difference matrix with unique
self-or-successor linkages.  Finally, the MDS parity-check theorem converts
selected columns in each local block into the capacity
$\min\{|U\cap J|,m-\sigma\}$.




\subsection{Formal verification audit and reproduction}
\label{app:lean-audit}

\paragraph{Trust boundary.}

The formalization uses no \texttt{sorry}, \texttt{admit}, or project-local
axiom.  A dedicated module applies \texttt{\#print axioms} to the theorems
stated in the paper.  Those theorems depend only on the standard Lean/mathlib
foundations.  The audit reports \texttt{propext},
\texttt{Classical.choice}, and \texttt{Quot.sound}.  Native decision procedures
are used only for labeled small-instance regression tests.  No Lean-formalized theorem stated
in the paper depends on a generated native-decision axiom.  The fixed-code
spectra and distances have kernel-checkable proofs.

The command \texttt{lake build BAA.AxiomAudit} checks
the complete development and prints the dependency audit.  Exhaustive
regressions cover the binary accumulator and difference-linkage conventions
through length eight and the nonbinary placement formula through length seven.
These tests supplement the symbolic all-length theorems; they do not replace
them.

The Lean development does not formalize the quadratic-residue constituents or
their published spectra.  It also does not claim the one-accumulator asymptotic
limit, \eqref{eq:B6}--\eqref{eq:B12}, \eqref{eq:6}--\eqref{eq:7},
\eqref{eq:L3}--\eqref{eq:L7}, any numerical transition, the $O(n)$
 cost model, or concentration of structural containers.  These claims require
 external coding-theory results, analytic estimates, certified optimization, a
 computational cost model, or new probabilistic arguments.  The RA$^t$
 parameter selection and configuration measurements are likewise outside the Lean development.
 Theorem~\ref{thm:rs-finite} additionally uses the analytic reductions in
 Appendix~\ref{app:rs-certificate} and the accompanying rational/Arb
 verifier~\cite{chosenbaaartifact}; those components are not Lean-formalized.
 The RA$^t$ finite sums have not received a complete numerical certificate.


